% profiles/arizona-state-university-campus-immersion.tex — Arizona State University
% ============================================================================
% source: https://graduate.asu.edu/current-students/completing-your-degree/formatting-your-thesis-or-dissertation/asu-graduate-college
%         (Graduate College, the page that serves the Format Manual)
% source: https://graduate.asu.edu/sites/default/files/2022-02/asu-graduate-college-format-manual.pdf
%         ("ASU Graduate College Format Manual", footed "February 2022" on all
%          twenty pages; the annotated title-page specimen is page 5 of 20)
% source: https://graduate.asu.edu/current-students/completing-your-degree/formatting-your-thesis-or-dissertation
%         (the Formatting page, which links the accessibility guidance below)
% source: https://graduate.asu.edu/sites/g/files/litvpz641/files/2026-08/2026%2008-27%20Best%20Practice%20Creating%20Accessible%20ETD%20v2%20%28HS%29.pdf
%         ("Best practice for students and faculty advisors: Creating an
%          accessible Electronic Thesis/Dissertation (ETD)", dated 2026-08-27 in
%          its own asset path)
% accessed: 2026-09-09
% checked:  texlib-thesis-profile-round
%
% VERIFIED and set here: geometry, spacing, the title page, and the absence of a
% separate approval page.
% NOT set:
%   * \thesissetchapteropening — no deeper chapter-opening margin is prescribed.
%   * \thesisrequiretagging, \thesisrequiredocumentlanguage,
%     \thesisrequirepdfstandard — see the accessibility note; ASU's guidance is
%     published as BEST PRACTICE, not as a filing requirement.
%
% NAMING. IPEDS records this unit as "Arizona State University Campus
% Immersion", an administrative name that appears nowhere in ASU's own
% documents. The Format Manual's title-page specimen prints "ARIZONA STATE
% UNIVERSITY", and that is what \thesisinstitution carries and what the page
% below sets.
%
% ============================================================================
% READ THIS BEFORE FILING: THIS CLASS DOES NOT SATISFY ASU'S TYPEFACE RULE.
% ============================================================================
% Unlike every other institution in this directory so far, ASU prescribes the
% FACE and the SIZE together, from a closed list:
%
%     "The Graduate College requires that students use one of the TrueType fonts
%      listed below. You should retain the same font and font size throughout
%      your document (preliminary, main text, back matter pages); the only
%      exception is endnotes and footnotes which may be in a smaller point size."
%
%         Arial 10pt · Century 11pt · Garamond 12pt · Georgia 11pt
%         Sans Serif 10pt · Tahoma 10pt · Times New Roman 12pt · Verdana 10pt
%
% texlib-thesis sets Latin Modern at 12pt. Latin Modern is on no part of that
% list, so a document built with the class as it stands is OUT OF SPEC at ASU —
% and this is the one requirement a profile has no field for. It cannot be
% fixed here without a profile reaching past the filing figures into the
% document's typography, which is not what a profile owns.
%
% The filer has to do it, in the document preamble, e.g.
%
%     \setmainfont{Times New Roman}      % 12pt, per the list above
%
% (fontspec is already loaded by the class, and lualatex will find the system
% face on Windows and macOS.) Note the manual's further warning: "There are
% multiple versions of Times New Roman font that vary slightly. If using Times
% New Roman, make sure that only one version is used consistently throughout
% entire document." A Times CLONE — TeX Gyre Termes, Nimbus Roman — is metric-
% compatible but is not the named face; whether the format reviewer accepts one
% is not something these sources answer, so do not assume it.
%
% ONE MORE THING THE PROFILE CANNOT EXPRESS. Pagination: "Title page is not
% paginated"; "All preliminary pages (title page to preface) are paginated with
% lowercase Roman numerals, starting with `i' on the abstract"; body pages in
% arabic "starting with `1' on the first page"; "Place all page numbers at the
% bottom of the page, centered". The abstract-is-page-i rule in particular is
% not something the class's \frontmatter produces, and the optional copyright
% page "does not change pagination". A filer must set that up.

\thesisinstitution{Arizona State University}

% "Every page of your document must meet the margin requirements of 1.25 inches
% on the left and right, and 1 inch on the top and bottom. All materials
% including appendices, if you choose to include them, must meet the margin
% requirements."
%
% ASU is the first institution in this directory whose margins are NOT square.
% The manual's annotated specimen labels the same four figures on the page
% itself (1.25" at each side, 1" top and bottom), so this is confirmed twice.
\thesissetgeometry{left=1.25in, right=1.25in, top=1in, bottom=1in}

% "All text must be double-spaced, except: block quotes, appendices,
% table/figure captions, material in tables, footnotes, endnotes, reference
% citations, and the optional biographical sketch. You must single-space
% individual footnotes and reference entries, then double-space between each
% note and entry." The title page has its own single-spaced blocks, handled
% below.
\thesissetspacing{double}

% --- Accessibility -----------------------------------------------------------
% ASU publishes real, specific accessibility guidance — and publishes it as
% GUIDANCE. The document is titled "Best practice for students and faculty
% advisors", and its own framing sentence is:
%
%     "This document provides guidance to help ensure that ETDs comply with Web
%      Content Accessibility Guidelines (WCAG) 2.1, level AA in order to adhere
%      to state and federal law, the ASU IT Accessibility Standard and ASU's
%      charter."
%
% The feature list beneath it is written in the declarative ("Text is selectable
% and searchable", "All images include alt text", "Include tags") — describing
% what an accessible document looks like, not stating what the Graduate College
% will reject a filing over. The February 2022 Format Manual, which IS the
% document a format review is run against, contains the word "accessib" nowhere
% at all; its only adjacent line is that four of the eight permitted faces "are
% designed for easy screen readability and are highly recommended".
%
% So the checkable declarations stay UNSET. That follows the reading already
% applied to Texas A&M: encouragement is not a requirement, and declaring one
% ASU never published would put words in the Graduate College's mouth. The one
% sentence that pulls the other way is the second entry below — an ASU standard
% is asserted to COVER ETDs — and it is quoted here rather than acted on so a
% reviewer can weigh it without re-reading the source.
\thesisaccessibilityrequirement{Published as best practice, not as a filing rule}
	{ASU's accessibility document is titled ``Best practice for students and
	 faculty advisors'' and says it ``provides guidance to help ensure that ETDs
	 comply with Web Content Accessibility Guidelines (WCAG) 2.1, level AA''.
	 The Format Manual that a format review is run against says nothing about
	 accessibility at all.}
\thesisaccessibilityrequirement{Scope of the ASU IT Accessibility Standard}
	{``ASU is committed to providing digital accessibility for resources used by
	 members of the university community and the general public. The ASU IT
	 Accessibility Standard covers electronic theses and dissertations (ETDs)
	 published on the ASU Digital Repository.''}
\thesisaccessibilityrequirement{Text readability}
	{``Text is selectable and searchable, not just an image. Fonts are embedded
	 (also a requirement for ProQuest publication). Avoid including text that
	 conveys important information within images.''}
\thesisaccessibilityrequirement{Proper headings}
	{``Use structured headings (H1, H2, etc.) for easy navigation. Avoid
	 skipping heading levels and keep formatting consistent.''}
\thesisaccessibilityrequirement{Alt text for images}
	{``All images include alt text describing the image for visually impaired
	 readers, or are marked as decorative. Tables may also include alt text
	 summaries when needed.''}
\thesisaccessibilityrequirement{Tagged content}
	{``Include tags that help assistive technology understand structure.''}
\thesisaccessibilityrequirement{Logical reading order}
	{``Content flows in a natural, logical order. Avoid a disorganized or
	 jumbled reading order. Use simple layouts for tables and include tagged
	 heading rows/columns.''}
\thesisaccessibilityrequirement{Document title in the file properties}
	{``Include the document title under `properties' in final PDF (separate from
	 the file name).''}

% --- The approval date ---------------------------------------------------------
% ASU's title page carries TWO dates and the class has one. The upper one is the
% defence: "Date: month and year of your oral defense; no comma". The lower one
% is the conferral month: "Date: month (either May, August, or December)", and
% that is \thesis@date. Set the defence date with \asuapproved{April 2027}.
% Unset, the page falls back to \thesis@date, which is very likely wrong — so it
% is worth setting even when the two months happen to agree.
\newcommand{\asu@approved}{}
\newcommand{\asuapproved}[1]{\renewcommand{\asu@approved}{#1}}

% --- Title page ----------------------------------------------------------------
% Reproduced from the Format Manual's annotated specimen (page 5 of 20). The
% specimen's own text, with the manual's spacing annotations applied:
%
%     Eurocentric Myths
%     in Three Quechua Novels          <- title, double spaced, Title Case
%     by                               <- all lowercase
%     Nigel St. John-Smith
%     [5 single line spaces]
%     A Dissertation Presented in Partial Fulfillment
%     of the Requirements for the Degree
%     Doctor of Philosophy             <- single spaced; degree name only
%     [10 single line spaces]
%     Approved April 2013 by the
%     Graduate Supervisory Committee:
%
%     Humbolt Volker, Chair
%     Jean-Pierre Fusquiere
%     Mercees Tomlin
%     ARIZONA STATE UNIVERSITY         <- ALL CAPITAL LETTERS
%     May 2013
%
% Every bullet the manual prints for this page is honoured:
%   - "Title: double-spaced, Title Case, center aligned" — so \@title is emitted
%     as written and never uppercased, and never bolded (nothing on this page is).
%   - "``by'': double-spaced; all lowercase letters, center aligned".
%   - "Name: what is included in the ASU system".
%   - "Clause ``A Thesis Presented in Partial Fulfillment'': 5 single line spaces
%     after name" — the specimen prints "A Dissertation Presented..." for a
%     doctorate, so the noun comes from \thesis@DocNoun.
%   - "Description: single-spaced; insert name of degree only".
%   - "Clause ``Approved [Month] [Year] by the'': ten single line spaces after
%     the degree".
%   - "Block 3: ``Graduate Supervisory Committee'' in title case followed by a
%     colon (:); insert names of committee members; comma and space before Chair
%     and Co-Chair" — so a member with no role prints bare, and the class's
%     signature rule must go: ASU collects no signatures on this page.
%   - "University: ALL CAPITAL LETTERS".
%   - "Title page is not paginated".
\renewcommand{\thesistitlepage}{%
	\loadgeometry{nopagenumber}%
	\begin{titlepage}%
		\thispagestyle{empty}\centering
		% The specimen's title block begins AT the 1in top margin -- no leading
		% fill. The single slack in the page is above the university line, which
		% the manual's own annotation places at the foot.
		{\doublespacing
		 \thesis@tagtitle{Title}{\@title}\par
		 by\par
		 \@author\par}%
		% "5 single line spaces" and, below, "ten single line spaces" -- SINGLE,
		% so they are measured against the single-spaced baseline, not against
		% the page's double-spaced one. Taking them at the ambient (double)
		% \baselineskip overruns the page and pushes the last committee member
		% and the university line off it.
		{\singlespacing
		 \vspace{5\baselineskip}
		 A \thesis@DocNoun\ Presented in Partial Fulfillment\\
		 of the Requirements for the Degree\\
		 \thesis@degree\par
		 \vspace{10\baselineskip}
		 Approved \ifx\asu@approved\@empty\thesis@date\else\asu@approved\fi\
		 by the\\
		 Graduate Supervisory Committee:\par
		 \vspace{\baselineskip}
		 % No rules: ASU prints typed names only, with a role label on the chair
		 % and co-chairs and on nobody else.
		 \begingroup
		 \renewcommand{\thesis@memberline}[2]{%
			##1\ifx\relax##2\relax\else, ##2\fi\par
		 }%
		 \thesis@committee
		 \endgroup}%
		\vfill
		{\doublespacing
		 ARIZONA STATE UNIVERSITY\par
		 \thesis@date\par}%
	\end{titlepage}%
	\loadgeometry{normal}%
}

% --- Committee approval page ---------------------------------------------------
% None. The Graduate Supervisory Committee is named on the title page above, and
% the manual's "Contents and Order" list runs Title Page, Copyright page,
% Abstract, Dedication, Acknowledgments, Table of Contents, ... with no approval
% or signature leaf anywhere in it.
\thesisnoapprovalpage
